\section{Conclusions}\label{sec:conclusions}

Scored on the spawning-stock biomass of Northern cod under a retention rule whose scoring core was coded before the first pass and applied unchanged, no member of a five-module surplus-production ladder improves on last-value persistence at one- and five-year horizons, on either the 1983--2015 NCAM series or the 1954--2024 xteNCAM series. The margins are not within the 5\% band on any specification or horizon; on the longer series they exclude zero. Decomposing the pooled score shows that the deficit is carried by the 1991--1995 collapse, which no member of the class can reproduce at the recorded catch because the collapse transitions removed more biomass than the catch record accounts for; that the autonomous modules forecast extinction at every post-collapse origin because training-mean catch exceeds the production of a depleted stock; and that the catch-supplied modules carry a linear low-biomass growth term with one-step skill on one specification that does not extrapolate to five years or to the other series. The lower equilibrium that makes the autonomous map decline is generated by the removals, not by depensation, and the depensation parameter is unidentified. Simulation under known truth makes the verdict strong against the autonomous module and weak against the others, and shows that retention, had it occurred, would not have identified which structure was present. After 2015 the reference point, the assessment vintage and the observing system moved more than any forecast did. For stock assessment the implication is procedural: claims that a structural element improves prediction should be made on the state, against persistence, at the horizons advice uses, under a rule declared in advance and accompanied by its operating characteristics; the present test supplies a template and a negative result for one stock on which such claims have been frequent.
